\lesson{Type members and abstraction}\label{sec:objects}
\focus{A type member separates an object's type definition from the bounds available to its clients.}

In Dependent Object Types (DOT), the source calculus family studied here,
an object can define a type member as well as hold value fields.
For a type $W$, consider the source-style sketch
\begin{lstlisting}
let x = { type A = W }
\end{lstlisting}
The selection $x.A$ names the type member of this particular object.
Another object $y$ may define a member with the same label $A$, but the
selections $x.A$ and $y.A$ are distinct names. Relating them requires
information about their definitions or bounds.

For the running example, fix
\[
  W=\top\to\top,\qquad f=\lambda(u:\top).u.
\]
Here $\top$ is the universal type, so $f:W$. We also use $\bot$ for the
empty type: a closed well-typed computation cannot return a value of
that type, although it may diverge. Using a function type for $W$ keeps
the example within the formal calculus, without additional primitives.

\subsection*{Definitions and interfaces}
The declaration $A=W$ defines the member $A$ as $W$. An interface
$\iface{S}{T}$ specifies the bounds available to a client:
\[
  x:\iface{S}{T}
  \quad\Longrightarrow\quad
  S<:x.A\quad\hbox{and}\quad x.A<:T.
\]
The lower bound permits a value of type $S$ to be used at the abstract
type $x.A$; the upper bound permits a value of type $x.A$ to be used
at $T$.

The exact interface $\iface{W}{W}$ provides both $W<:x.A$ and $x.A<:W$.
The less informative interface $\iface{\bot}{\top}$ provides only
$\bot<:x.A$ and $x.A<:\top$. The object can have the same definition in
both cases; widening the interface restricts the facts available to the
client.

This is the abstraction provided by a module that exposes a type and
operations on it while hiding its representation. The running object
contains only a type definition, allowing us to study the reasoning
about type members before introducing fields and methods.

\subsection*{The running program}
Figure~\ref{fig:program} gives the running program and its three subtyping
obligations. Subsequent sections translate its derivation, construct the
corresponding store, and execute the final application.


\begin{figure}[!htbp]
\begin{examplepanel}{srccol}{Source program \(P\) --- the running example}
\src{\(W=\top\to\top\). Let annotations specify the assigned types.}
\begin{lstlisting}[style=source]
let x         = { type A = W }
let f         = lambda (u: Top). u
let hidden    : x.A = f          -- L: W <: x.A
let recovered : W   = hidden     -- H: x.A <: W
recovered f                     -- T: W <: Top
\end{lstlisting}
\end{examplepanel}
\begin{examplepanel}{srccol}{Subtyping obligations}
\[
\src{
 \begin{array}{ccl}
 \text{L: hide} & W<:x.A & \text{the lower bound of }x.A,\\
 \text{H: recover} & x.A<:W & \text{the upper bound of }x.A,\\
 \text{T: pass the argument} & W<:\top & \text{inclusion into the universal type}.
 \end{array}}
\]
\end{examplepanel}
\caption{The running source program, with the object viewed at
\(\iface{W}{W}\). The aliases present the identity first at \(x.A\)
and then at \(W\); the final application presents its argument at
\(\top\). This sketch exposes the member bounds directly.
Section~\ref{sec:evidence} translates its client body;
Section~\ref{sec:closed-program} gives the complete translation, including
the recursive opening required by the formal source.}
\label{fig:program}
\end{figure}



The binding of $\ty{hidden}$ uses the lower bound $W<:x.A$.
The binding of $\ty{recovered}$ uses the upper bound $x.A<:W$.
The final application uses $W<:\top$ to type its argument. Since both
$\ty{recovered}$ and $f$ denote the identity, evaluation returns the
identity function supplied as the argument.

The displayed program is an explanatory sketch, rather than accepted
Scala syntax. The formal source names intermediate expressions and
explicitly records the step that exposes an object literal's member
interface; these details are introduced where the translation needs them.

\remember{The stored definition $A=W$, the stable name $x.A$, and the
interface supplying facts about that name have distinct roles. An
interface change preserves the definition and the name.}

\exercise{bounds}{At the abstract interface $\iface{\bot}{\top}$, can a
client justify assigning our $f:W$ to a variable of type $x.A$ using the
bounds alone? Write the inclusion it would need and compare it with
the available lower bound.}

\lesson{From typing derivations to target terms}\label{sec:evidence}
\focus{The translation represents subtyping derivations as evidence and subsumption as a cast.}

The source subsumption rule combines a typing derivation $h$ with a
subtyping derivation $d$:
\[
\src{\frac{h:\Gamma\vdash u:S\qquad d:\Gamma\vdash S<:T}
           {\textsc{Sub}(h,d):\Gamma\vdash u:T}}.
\]
Subsumption changes the assigned type while preserving the source
expression $u$. The chosen conversion is therefore recorded in the
derivation rather than the term.

Write $\tr{h}$ for the target term translated from a source typing
derivation $h$, and $\tr{d}$ for the evidence translated from a source
subtyping derivation $d$. The subsumption clause is
\[
\tgt{\tr{\textsc{Sub}(h,d)}=\cast{\tr{h}}{\tr{d}}.}
\]
This translation clause recursively translates the two premises and
combines their results with the target $\ev{cast}$ constructor.

The target is called \FC{}. An expression $e:S\le T$ is called
\emph{inclusion evidence}, or an \emph{inclusion coercion}.
The target has no implicit subsumption: a cast records a type change
justified by evidence, with no potentially failing runtime type test.
Its typing rule checks the two translated premises:
\[
\tgt{\frac{\Gamma'\vdash \tr{h}:S'\qquad
           \Gamma'\vdash\tr{d}:S'\le T'}
          {\Gamma'\vdash\cast{\tr{h}}{\tr{d}}:T'}}.
\]
Primes denote the translated context and types. We use the same notation
for $W$ and $x.A$ on both sides of the translation.

A target \emph{atom} is a variable, possibly surrounded by typing
wrappers such as casts. Its \emph{root} is the underlying variable:
both $f$ and $\cast{f}{e}$ have root $f$. Atoms serve as the receivers
and arguments of operations.

The \emph{erasure} $\erase{t}$ removes typing annotations and evidence
wrappers, yielding the underlying program. In particular,
$\erase{\cast{u}{e}}=\erase{u}$.
Erasure relates program syntax; the typed target's operational semantics
may still consult evidence.

\subsection*{Translating the two bound selections}
Consider an occurrence of $f$ under the parameter context
$\src{\Gamma=x:\iface{W}{W},\,f:W}$.
This is the client part of the example: for now, assume that its object
parameter already exposes the exact bounds. We will translate this body
first. Section~\ref{sec:closed-program} then supplies the literal's
introduction and recursive opening, displays the complete target program,
and substitutes the allocated atoms into this same body. No assumption
about a stored definition is needed for the present calculation.
Reading the assumed interface gives two source derivations:
\[
\src{\begin{aligned}
d_\ell&:\Gamma\vdash W<:x.A,\\
d_h&:\Gamma\vdash x.A<:W.
\end{aligned}}
\]
Write the translated evidence using the abbreviations
\[
 \ell_x=\ev{lower}_A(x):W\le x.A,\qquad
 h_x=\ev{upper}_A(x):x.A\le W.
\]
Section~\ref{sec:telescopes} expands each abbreviation into an indexed
member read. Figure~\ref{fig:derivation} translates the two subsumption
steps from their premises: variable typing yields the atom $f$ and the
lower-bound derivation yields $\ell_x$, which the inner \textsc{Sub}
combines into $\cast{f}{\ell_x}$. The upper-bound derivation then yields
$h_x$, and the outer \textsc{Sub} adds the second cast.


\begin{figure}[!htbp]
\begin{examplepanel}{srccol}{Source derivation --- two subsumptions on the same occurrence of \(f\)}
\small
\src{\(\Gamma=x:\iface{W}{W},\,f:W\).
Here \(h_f\) derives variable typing for \(f\), and \(d_\ell,d_h\)
derive the lower and upper bounds of \(x.A\).}
\[
\src{
\inferrule*[right=\textsc{Sub}]{
 \inferrule*[right=\textsc{Sub}]{
 h_f:\Gamma\vdash f:W\\
 d_\ell:\Gamma\vdash W<:x.A
 }{\Gamma\vdash f:x.A}\\
 d_h:\Gamma\vdash x.A<:W
}{\Gamma\vdash f:W}}
\]
\end{examplepanel}
\begin{examplepanel}{tgtcol}{Translation of the derivation}
\[
\tgt{
\begin{aligned}
 \tr{h_f}&=f,\\
 \tr{d_\ell}&=\ell_x=\ev{lower}_A(x):W\le x.A,\\
 \tr{d_h}&=h_x=\ev{upper}_A(x):x.A\le W,\\
 \tr{\textsc{Sub}(h_f,d_\ell)}&=\cast{f}{\ell_x},\\
 \tr{\textsc{Sub}(\textsc{Sub}(h_f,d_\ell),d_h)}
   &=\cast{\cast{f}{\ell_x}}{h_x}.
\end{aligned}}
\]
\end{examplepanel}
\begin{examplepanel}{runcol}{Erasure of the target term}
\[
\run{\erase{\cast{\cast{f}{\ell_x}}{h_x}}=f.}
\]
\end{examplepanel}
\caption{Each \textsc{Sub} node contributes a target \(\ev{cast}\),
while the source expression remains \(f\) throughout the derivation.
The calculation assumes parameter binders for \(x\) and \(f\).
Figure~\ref{fig:let-translation} extends it to the program's let bindings.}
\label{fig:derivation}
\end{figure}



\subsection*{Dependence on the source derivation}
The judgment $\src{\Gamma\vdash f:W}$ admits several derivations.
Variable typing translates to $f$, while the two subsumptions translate
to $\cast{\cast{f}{\ell_x}}{h_x}$. A third derivation first composes
the bounds:
\[
\src{\frac{d_\ell:\Gamma\vdash W<:x.A\qquad
           d_h:\Gamma\vdash x.A<:W}
          {\textsc{Trans}(d_\ell,d_h):\Gamma\vdash W<:W}}.
\]
Source transitivity becomes an evidence constructor:
\[
\tgt{\tr{\textsc{Trans}(d_\ell,d_h)}
       =\ev{trans}(\ell_x,h_x).}
\]
One subsequent subsumption therefore gives
$\cast{f}{\ev{trans}(\ell_x,h_x)}$.
The three target expressions are syntactically distinct, although all
erase to $f$. In particular, nested casts and a cast by composed evidence
are different constructor expressions. The translation is thus a
function of the source derivation, not only of its concluding term.

We abbreviate $\ev{trans}(p,q)$ by $p;q$, in derivation order:
\[
 \frac{p:S\le M\qquad q:M\le T}{p;q:S\le T}.
\]
Composition requires matching intermediate endpoints. Here the chain is
$W\xrightarrow{\ell_x}x.A\xrightarrow{h_x}W$.

\subsection*{Let translation and evaluation}
The translation preserves the running program's let bindings:
\[
\tgt{
\tr{\textsc{Let}(h_1,h_2)}
=\ev{let}\ y=\tr{h_1}\ \ev{in}\ \tr{h_2}.}
\]
Figure~\ref{fig:let-translation} applies this clause to the body of the
running example. The first right-hand side uses the lower-bound cast.
The second uses the upper-bound cast on the newly bound name
$\ty{hidden}$. The application also needs the argument cast
$\ev{top}(W):W\le\top$.

Application translates the variable typing derivations of its operands
to atoms. Subsumption in such a derivation contributes a cast wrapper,
so the source argument $f$ becomes $\cast{f}{\ev{top}(W)}$.
The full language distinguishes casts of terms, values, and atoms;
we use the same notation when the syntactic position determines the sort.


\begin{figure}[!htbp]
\begin{examplepanel}{srccol}{Source body --- after binding \(x\) and \(f\)}
\begin{lstlisting}[style=source]
let hidden    : x.A = f in
let recovered : W   = hidden in
recovered f
\end{lstlisting}
\end{examplepanel}
\begin{examplepanel}{tgtcol}{Target translation with explicit casts}
\tgt{\(\ell_x=\ev{lower}_A(x)\), \(h_x=\ev{upper}_A(x)\),
\(t_W=\ev{top}(W):W\le\top\).}
\begin{lstlisting}[style=target,mathescape=true]
let hidden    = cast(f, $\ell_x$) in
let recovered = cast(hidden, $h_x$) in
recovered (cast(f, $t_W$))
\end{lstlisting}
\end{examplepanel}
\begin{examplepanel}{tgtcol}{Evaluation by substitution of the aliases}
\small
\[
\tgt{\begin{aligned}
 &\ev{let}\ \ty{recovered}=\cast{\cast{f}{\ell_x}}{h_x}\ \ev{in}\
       \ty{recovered}\ \cast{f}{t_W}\\
 &\hspace{5em}\longrightarrow^{*}
       \cast{\cast{f}{\ell_x}}{h_x}\ \cast{f}{t_W}.
\end{aligned}}
\]
The first binding substitutes \(\cast{f}{\ell_x}\) for \(\ty{hidden}\);
the second substitutes the nested cast for \(\ty{recovered}\).
The relation \(\longrightarrow^{*}\) includes the administrative steps
that rearrange pending casts.
\end{examplepanel}
\caption{Translation preserves the source let structure; evaluation
produces nested casts by substituting the let-bound atoms. The resulting
application erases to \(f\,f\). As in Figure~\ref{fig:derivation}, the
exact interface is assumed directly.
Section~\ref{sec:closed-program} adds the literal's introduction and
recursive opening, then instantiates its binders with atoms over the store.}
\label{fig:let-translation}
\end{figure}



During evaluation, substitution of the let-bound atoms produces the
nested wrappers. This display still has the parameters $x$ and $f$ free.
Closing them is a further calculation, given in
Section~\ref{sec:closed-program}; it must justify both the object
interface used by the evidence and the function supplied to the call.

\subsection*{Checking evidence endpoints}
The expression $\ell_x;\ev{top}(W)$ fails the composition rule:
its first component ends at $x.A$ and its second starts at $W$.
Inserting the upper-bound evidence gives the well-typed composition
$\ell_x;h_x;\ev{top}(W)$. Alternatively,
$\ell_x;\ev{top}(x.A)$ reaches $\top$ directly from $x.A$.

Even if a later store resolves $x.A$ to $W$, checking still compares
syntactic endpoints. Store-based normalization establishes what the
evidence means at runtime.

\subsection*{Erasure and operational correspondence}
The erasure equation $\erase{\cast{u}{e}}=\erase{u}$ helps show that
translation preserves the underlying erased program.
The typed target machine still consults a function atom's
evidence to perform a call. A simulation theorem relates those typed
steps to the erased program's execution.

\remember{Translation constructs evidence from source derivations;
evaluation substitutes atoms and consults their evidence; erasure
removes the evidence to relate the underlying programs.}

\exercise{order}{What does $h_x;\ell_x$ prove? Is it the same expression
as $\ell_x;h_x$? Also give explicit evidence for $W\le\top$ that goes
through $x.A$.}

\exercise{translate}{The source occurrence $f:W$ is first used at
$x.A$, then at $\top$ by the source \textsc{Top} rule $x.A<:\top$.
Give its target term when there are two \textsc{Sub} nodes, and when
there is one \textsc{Sub} over a \textsc{Trans} node. What does each erase to?}


\lesson{Interfaces as telescopes}\label{sec:telescopes}
\focus{The target represents an object interface as indexed propositions under a shared self binder.}

The interface $\iface{W}{W}$ specifies two bounds on its object's type
member. Its target representation binds a self variable $s$ in those
propositions:
\[
 \ty{Exact}=\obj{W\le s.A,\ s.A\le W}.
\]
This type describes an object satisfying the listed propositions, with
$s$ denoting the object itself. Here $\mu$ binds self in a proposition
list; this example does not involve recursive unfolding into a larger
type structure.

The development calls this list a \emph{telescope}. Its propositions are
indexed from zero:
\[
 \begin{array}{ccl}
 \text{position }0&:&W\le s.A,\\
 \text{position }1&:&s.A\le W.
 \end{array}
\]
All propositions share the self binder; individual propositions do not
bind term variables. Indexed selection therefore represents both lower
and upper bounds without a dedicated constructor for either.

\subsection*{Selecting a proposition}
For an atom $a:\ty{Exact}$, inclusion evidence is obtained with
\[
 \ev{member}\ a\ (\ev{refl}\ \ty{Exact})\ i.
\]
The atom $a$ identifies the object, the coercion selects an interface
reachable from its current type, and the index $i$ selects a proposition
in that interface. Since $a$ already has type $\ty{Exact}$, the coercion
is reflexivity, $\ev{refl}\ \ty{Exact}:\ty{Exact}\le\ty{Exact}$.

If $a$ is rooted at $x$, reading position 1 replaces $s$ by $x$:
\[
 \ev{member}\ a\ (\ev{refl}\ \ty{Exact})\ 1:x.A\le W.
\]
This expands the abbreviation $\ev{upper}_A(a)$. Position 0 similarly
gives $\ev{lower}_A(a):W\le x.A$. The abbreviations suppress only the
index and reflexivity evidence.

\subsection*{Roots and dependent substitution}
For the current fragment, atoms and their roots are defined by
\[
 a::=x\mid\cast{a}{e},\qquad
 \ev{root}(x)=x,\quad
 \ev{root}(\cast{a}{e})=\ev{root}(a).
\]
Suppose $a=\cast{x}{c}:\ty{Exact}$. The selected type is still $x.A$.
There is no new type name $(\cast{x}{c}).A$ in the type syntax.

When a function body checked under
$y:\ty{Exact}$ receives the atom $\cast{x}{c}$, occurrences of
$y$ as an atom become $\cast{x}{c}$, while occurrences of $y.A$ in
types become $x.A$. The wrapper supplies the required interface; the
root supplies the name.

Preserving the root ensures that the lower- and upper-bound proofs
continue to refer to the same selection after an interface change.
Their shared endpoint $x.A$ is what permits their composition in the
running example.

\subsection*{Equality propositions}
The literal that defines $A=W$ has the more precise type
\[
 \ty{Lit}=\obj{s.A\equiv W}.
\]
Its single equality proposition yields the two inclusions by reading
the equality in opposite directions. The fully abstract interface is
\[
 \ty{Abs}=\obj{\bot\le s.A,\ s.A\le\top}.
\]
The next two sections construct and evaluate the conversions
$\ty{Lit}\le\ty{Exact}\le\ty{Abs}$.

\exercise{root}{Let $a=\cast{\cast{x}{c_1}}{c_2}:\ty{Abs}$.
What are the endpoints obtained by reading positions 0 and 1?
Which part of the answer depends on the casts, and which part does not?}

\lesson{Constructing interface coercions}\label{sec:recipes}
\focus{An object coercion derives each destination proposition from the source telescope.}

The two interface conversions require evidence
\[
 \ev{expose}:\ty{Lit}\le\ty{Exact},\qquad
 \ev{forget}:\ty{Exact}\le\ty{Abs}.
\]
The \emph{source} and \emph{destination} of a coercion are its input and
output interfaces, both within the target language. For each destination
proposition, we construct a recipe deriving it from a source proposition.

\subsection*{From a definition to exact bounds}
The source telescope of $\ty{Lit}$ contains one slot:
$q_0:s.A\equiv W$. The destination $\ty{Exact}$ requires two facts:
\[
 \begin{array}{lll}
 \text{required lower bound}&W\le s.A&
       \text{read }q_0\text{ backwards},\\
 \text{required upper bound}&s.A\le W&
       \text{read }q_0\text{ forwards}.
 \end{array}
\]
Equality supports both inclusion directions; inclusion evidence alone
does not support reversal. Both output recipes may therefore use the
same input slot.

The name $q_0$ is a \emph{placeholder} for a source proposition, not a
fact already selected from a particular object. These recipes define a
conversion independently of the object to which it will be applied.
Their ordered list is called a \emph{morphism}.

\subsection*{From exact bounds to abstract bounds}
Now the source slots are
$p_0:W\le s.A$ and $p_1:s.A\le W$.
For the destination lower bound, prepend the inclusion from $\bot$:
\[
 \bot\xrightarrow{\ev{bot}(W)}W
      \xrightarrow{p_0}s.A.
\]
For the destination upper bound, append the inclusion into $\top$:
\[
 s.A\xrightarrow{p_1}W
      \xrightarrow{\ev{top}(W)}\top.
\]
These two recipes discharge the destination obligations using the
source's lower and upper slots, respectively.

\subsection*{Templates and morphisms}
Each of these recipes is represented by a \emph{template} with the three-part form
\[
 \boxed{\text{left side}\quad|\quad\text{source slot}\quad|\quad\text{right side}}.
\]
A side may be absent when its endpoint requires no further conversion.
The central slot is a \emph{hole}, instantiated by a fact about the
object when the coercion is applied. Figure~\ref{fig:recipes} gives all
four recipes and the source subtyping rule translated by $\ev{forget}$.

\begin{figure}[!htbp]
\begin{examplepanel}{srccol}{Source widening --- the declaration subtyping rule}
\[
\src{\inferrule*[right=\textsc{Typ}]{
 \bot<:W\\ W<:\top
}{\iface{W}{W}<:\iface{\bot}{\top}}}
\]
Widening decreases the lower bound and increases the upper bound.
\end{examplepanel}
\begin{examplepanel}{tgtcol}{Target interfaces}
\[
\tgt{\begin{aligned}
 \ty{Lit}&=\obj{s.A\equiv W},\\
 \ty{Exact}&=\obj{W\le s.A,\ s.A\le W},\\
 \ty{Abs}&=\obj{\bot\le s.A,\ s.A\le\top}.
\end{aligned}}
\]
\end{examplepanel}
\begin{examplepanel}{tgtcol}{Coercion templates --- one per destination proposition}
\small
\renewcommand{\arraystretch}{1.3}
\begin{tabularx}{\linewidth}{@{}l l >{\raggedright\arraybackslash}X l@{}}
\toprule
Conversion and output & Left side & Input hole & Right side\\
\midrule
$\ev{expose}$, lower & absent & equality 0, reversed & absent\\
$\ev{expose}$, upper & absent & equality 0, forward & absent\\
$\ev{forget}$, lower & $\ev{bot}(W)$ & inclusion 0 & absent\\
$\ev{forget}$, upper & absent & inclusion 1 & $\ev{top}(W)$\\
\bottomrule
\end{tabularx}
\end{examplepanel}
\caption{The coercion \(\ev{expose}:\ty{Lit}\le\ty{Exact}\) derives
exact bounds from the definition equality, while
\(\ev{forget}:\ty{Exact}\le\ty{Abs}\) translates the displayed
widening derivation. Each hole indexes the source telescope of its
coercion. The templates are independent of the object to which the
coercion will be applied.}
\label{fig:recipes}
\end{figure}


The indices always refer to the \emph{source} telescope of that conversion.
In particular, the two occurrences of index 0 in the first two rows
refer to the same equality.

In compact notation, omitting the constructors that attach a recipe to
its surrounding list, the upper recipe for $\ev{expose}$ is
\[
 \ev{le}\ \ev{none}\ (\ev{eq}\ 0)\ \ev{none}.
\]
The constructor $\ev{le}$ specifies an inclusion output, the two
$\ev{none}$ arguments specify absent sides, and $\ev{eq}\ 0$ reads
equality slot 0 forwards. Replacing $\ev{eq}$ by $\ev{eqSym}$ gives the
lower recipe.

The complete object coercion has the form
$\ev{obj}\ \ty{sourceTelescope}\ m$, where $m$ is the list of recipes.
It carries the source telescope. Its destination is determined by the
facts the recipes establish. The implementation represents these lists
with constructors; the bracket notation here displays their order.

\subsection*{Scope of the template sides}
In a basic template, present sides are scoped outside the self binder
and cannot mention it. For example, $\ev{bot}(W)$ refers only to $W$.
The recipe accesses $s$ through its hole.

The restriction concerns the scope of $s$; an outer variable may still
denote the same runtime object. An absent side can preserve an endpoint
such as $s.A$ without supplying evidence outside the binder.
As Section~\ref{sec:normalization} shows, this representation supports
composition by substitution into holes, without additional member
lookups through self.

\exercise{widen}{Build the conversion from $\ty{Exact}$ to
$\obj{\bot\le s.A,\ s.A\le W}$. Give the lower and upper recipes
in three-part form, and say which side of $\ev{forget}$ you removed.}

\lesson{Stores and views}\label{sec:store}
\focus{A typed store supplies the concrete definitions used to interpret interface evidence.}

The preceding interfaces and coercions can be checked before execution.
To interpret them at runtime, suppose the machine has allocated the
literal at $\ell$:
\[
 \sigma(\ell)=\{A=W\}.
\]
The store records the value itself. Write $\Gamma_\sigma$ for its typing
context: it records the precise type $\ty{Lit}$ and the definition of
$\ell.A$. A binder carrying definitions is called \emph{transparent}.
A function parameter
$y:\ty{Exact}$ is \emph{opaque}: its interface supplies assumptions, but
there is no stored definition for its $y.A$ while checking the body.

The transparent location admits definition evidence
$\ev{def}(\ell,A):\ell.A\equiv W$.
An opaque parameter supplies no definition to justify this constructor.
The coercion $\ev{expose}$ instead uses the equality in its source
telescope, so its construction is independent of the stored definition.

\subsection*{Precise allocation}
Suppose a let binds a literal already cast to $\ty{Exact}$:
\[
 \ev{let}\ y=\cast{\{A=W\}}{\ev{expose}}\ \ev{in}\ t.
\]
The machine stores the bare literal at a fresh location $\ell$ and
substitutes the atom $\cast{\ell}{\ev{expose}}$ for $y$ in the body.
The store therefore retains $\ty{Lit}$ and the definition $A=W$;
the client sees $\ty{Exact}$ through its evidence.

This allocation rule is part of the store invariant. Recording arbitrary
interfaces as store facts would admit assertions unsupported by the
stored values. Instead, the store retains precise definitions, and cast
evidence justifies each change of presentation.

\subsection{Closing the complete source translation}\label{sec:closed-program}
We can now connect the client calculation to the entire running program.
Use $x,f$ for its source binders and reserve $\ell,z$ for the two
locations allocated during execution. This keeps a parameter's assumed
interface distinct from a location's precise stored type.

\textbf{First supply the source object's interface.}
Let $D=\iface{W}{W}$. The source definition rule directly assigns
$A=W$ the exact declaration $D$. The object rule checks that definition
under a self binding at $\mu(s.D)$ that also records the literal's
definitions, and gives the literal type $\mu(s.D)$.
Here $W$ is independent of self, so checking the definition uses no
facts from that binding. After the first let, variable typing gives
$x:\mu(s.D)$, whereas reading a bound needs $x:D[x/s]$.
The source recursive-elimination rule supplies precisely that change:
\[
 \src{\frac{\Gamma\vdash x:\mu(s.D)}
                 {\Gamma\vdash x:D[x/s]}\ \textsc{Rec-E}.}
\]
It changes the derivation of the variable occurrence, not the source
expression. In this example $D$ contains no free occurrence of its
bound self, so $D[x/s]=D$. Nevertheless the rule is present in the
derivation and has a visible translation.

\textbf{Translate introduction and opening.}
The target literal has precise type $\ty{Lit}$, while the translation
of the source literal's assigned type is $\ty{Exact}$. The object
translation therefore inserts the coercion already constructed:
\[
 \tr{d_{\mathrm{obj}}}
       =\cast{\{A=W\}}{\ev{expose}}.
\]
For recursive opening, write
$P_E=[W\le s.A,\ s.A\le W]$ for the telescope of $\ty{Exact}$ and define
the exact wrapper abbreviation
\[
 R(c)=\ev{foldSelf}_{P_E}(\ev{unfoldSelf}(c)).
\]
The unfold wrapper opens the recursive presentation at the atom's root;
the fold wrapper packages the opened declaration's telescope as an
object interface. These are atom constructors, not evidence composition.
Here opening leaves the telescope's external types unchanged, so
$R(c):\ty{Exact}$ whenever $c:\ty{Exact}$.
Both constructors preserve the root. A view records evaluated interface
facts, as the next calculation will explain; their view rules give
$\view{R(c)}=\view{c}$ when interpreted over a typed store.
They are still present in the syntax and consume normalizer fuel.

Thus the bound derivations for this literal-origin source program
translate to $\ev{lower}_A(R(x))$ and $\ev{upper}_A(R(x))$.
The earlier parameter calculation used $\ev{lower}_A(x)$ and
$\ev{upper}_A(x)$ because it assumed $x:D$ directly.
This is the exact difference between the two derivations.

\textbf{Assemble the whole target.}
The identity lambda translates to itself: its body uses variable typing.
Each of the four lets uses the let clause from
Section~\ref{sec:evidence}. With the chosen lower-bound, upper-bound,
and argument subsumptions, the complete translation is
\[
\tgt{\begin{aligned}
 \ev{let}\ x&=\cast{\{A=W\}}{\ev{expose}}\ \ev{in}\\
 \ev{let}\ f&=\lambda(u:\top).u\ \ev{in}\\
 \ev{let}\ \ty{hidden}&=\cast{f}{\ev{lower}_A(R(x))}\ \ev{in}\\
 \ev{let}\ \ty{recovered}
   &=\cast{\ty{hidden}}{\ev{upper}_A(R(x))}\ \ev{in}\\
 &\ty{recovered}\ \cast{f}{\ev{top}(W)}.
\end{aligned}}
\]
Every constructor introduced by this derivation is now accounted for:
the object rule inserts $\ev{expose}$, \textsc{Rec-E} inserts $R$,
the three subsumptions insert casts, and lambda and let translation
retain the program structure. The notation still suppresses the
implementation's variable indices; $R$ and the member-read abbreviations
have explicit definitions.

\textbf{Close the body by allocation and substitution.}
After the first two bindings have allocated their bare values, the store
contains
\[
 \sigma(\ell)=\{A=W\},\qquad
 \sigma(z)=\lambda(u:\top).u,
 \qquad \Gamma_\sigma(\ell)=\ty{Lit},\quad
        \Gamma_\sigma(z)=W.
\]
The literal's wrapper is retained in the atom substituted for $x$:
\[
 a=\cast{\ell}{\ev{expose}},\qquad a^\circ=R(a),\qquad
 \theta=[x:=a,\ f:=z].
\]
For the remaining body $B$ of the displayed target, substitution gives
\[
\tgt{\begin{aligned}
 B\theta={}&
 \ev{let}\ \ty{hidden}=\cast{z}{\ev{lower}_A(a^\circ)}\ \ev{in}\\
 &\ev{let}\ \ty{recovered}
     =\cast{\ty{hidden}}{\ev{upper}_A(a^\circ)}\ \ev{in}\\
 &\ty{recovered}\ \cast{z}{\ev{top}(W)}.
\end{aligned}}
\]
In evidence receivers, $R(x)$ has become the whole atom $R(a)$.
In type selections, $x.A$ has become $\ell.A$, since
$\ev{root}(a)=\ell$. Substituting the two aliases next gives the
closed application
\[
 \cast{\cast{z}{\ev{lower}_A(a^\circ)}}{\ev{upper}_A(a^\circ)}
 \quad \cast{z}{\ev{top}(W)}.
\]
This is the application evaluated in Section~\ref{sec:application}.
The following view calculation supplies its remaining obligation:
the two member reads really normalize to identity at this store.
The companion's \texttt{FullSource} examples check the full source
derivation and its exact target translation; the machine calculation here
uses the stated allocation and substitution rules.


\subsection*{Concrete views}
A \emph{view} records the evaluated information at an atom's object
interface, in telescope order with one entry per proposition.
An inclusion entry contains a \emph{head form}, which records the outer
structure of the inclusion after following stored definitions.
We mark evaluated results with hats. The form $\form{id}$ says the
endpoints agree; $\form{bot}$ says the source endpoint is $\bot$;
$\form{top}$ says the destination endpoint is $\top$.
Thus $\ev{top}(W)$ is evidence to evaluate, whereas $\form{top}$ is its
evaluated form. Section~\ref{sec:normalization} adds function and object
forms and explains their computation.

An equality entry instead contains an equality marker. For our literal:
\[
 \view{\ell}=[\,\form{eq}\,].
\]
The marker $\form{eq}$ says that the equality's endpoints agree after
following the definitions. It does not record a new equality axiom.
The typing of the view relates that marker to the proposition
$\ell.A\equiv W$.

Let
\[
 a=\cast{\ell}{\ev{expose}}:\ty{Exact},\qquad
 b=\cast{a}{\ev{forget}}:\ty{Abs}.
\]

The full translated program uses the reopened receiver
$a^\circ=R(a)$ from Section~\ref{sec:closed-program}. Its wrappers
preserve the view as well as the root:
$\view{a^\circ}=\view{a}$ and $\ev{root}(a^\circ)=\ell$.
We first compute the view of $a$, then use this equality when reading the
program's bounds. Reading those bounds directly from the bare location
$\ell:\ty{Lit}$ would select from the wrong telescope.

To compute $\view{a}$, apply the two $\ev{expose}$ recipes to
$[\,\form{eq}\,]$. Each recipe reads the same equality slot. Forward or
backward, its endpoints resolve to the same type $W$, so each produces
the identity inclusion form:
\[
 \view{a}=[\,\ev{le}\,\form{id},\ \ev{le}\,\form{id}\,].
\]
Here $\ev{le}$ tags an inclusion entry, distinguishing it from an
equality entry.

Applying $\ev{forget}$ uses composition of head forms, written $F;G$.
This operation computes with normalized data, whereas $p;q$ constructs
an evidence expression. The first recipe prepends $\form{bot}$ to the
slot's $\form{id}$; the second appends $\form{top}$:
\[
 \begin{aligned}
 \text{lower output}&:\ \form{bot};\form{id}=\form{bot},\\
 \text{upper output}&:\ \form{id};\form{top}=\form{top}.
 \end{aligned}
\]
Thus
\[
 \view{b}=[\,\ev{le}\,\form{bot},\ \ev{le}\,\form{top}\,].
\]
The same literal now supplies the weaker facts of $\ty{Abs}$.
From that interface alone, the client cannot recover $W$;
the store still records that definition. Figure~\ref{fig:presentations}
records the three views and the normalization of the running program's
member evidence. There $\hnf{e}$ denotes head normalization of evidence,
the computation that produces its head form.

\needspace{4\baselineskip}
The notation $\operatorname{resolve}_{\Gamma_\sigma}(T)$ means following
the definitions in $\Gamma_\sigma$ to expose the outer shape of $T$.
Here it turns $\ell.A$ into $W$.


\begin{figure}[!htbp]
\begin{examplepanel}{runcol}{Stored definition and type resolution}
\[
 \run{\sigma(\ell)=\{A=W\},\qquad
       \operatorname{resolve}_{\Gamma_\sigma}(\ell.A)=W.}
\]
\end{examplepanel}
\begin{examplepanel}{tgtcol}{Three target atoms and their concrete views}
\small
\renewcommand{\arraystretch}{1.35}
\begin{tabularx}{\linewidth}{@{}l >{\raggedright\arraybackslash}X l@{}}
\toprule
Interface & Atom rooted at \(\ell\) & View\\
\midrule
$\ty{Lit}$ & $\ell$ & $[\,\form{eq}\,]$\\
$\ty{Exact}$ & $a=\cast{\ell}{\ev{expose}}$
 & $[\,\ev{le}\,\form{id},\,\ev{le}\,\form{id}\,]$\\
$\ty{Abs}$ & $b=\cast{a}{\ev{forget}}$
 & $[\,\ev{le}\,\form{bot},\,\ev{le}\,\form{top}\,]$\\
\bottomrule
\end{tabularx}
\end{examplepanel}
\begin{examplepanel}{tgtcol}{Normalization of the running program's member evidence}
\[
\tgt{\begin{aligned}
 \ev{lower}_A(a^\circ)&:W\le \ell.A
     &\hnf{\ev{lower}_A(a^\circ)}&=\form{id},\\
 \ev{upper}_A(a^\circ)&:\ell.A\le W
     &\hnf{\ev{upper}_A(a^\circ)}&=\form{id}.
\end{aligned}}
\]
\end{examplepanel}
\caption{Applying coercion templates to the stored object's facts
produces concrete views. The atom \(a:\ty{Exact}\) instantiates the
allocated client binding, and $a^\circ=R(a)$ is its reopened
member receiver. Their root \(\ell:\ty{Lit}\) supplies the dependent
name \(\ell.A\). The abstract presentation exposes weaker interface facts
while the store retains the definition.}
\label{fig:presentations}
\end{figure}



\remember{Evidence expressions, normalized object recipes, and concrete
views are distinct representations. Normalized recipes retain holes;
applying them to a stored object's facts produces a concrete view.}

\exercise{views}{Through $a$, what head form does reading the upper bound
produce? What about through $b$? Explain why these answers can differ
even though $\ev{root}(a)=\ev{root}(b)=\ell$.}

\lesson{Evidence normalization and composition}\label{sec:normalization}
\focus{Head forms support composition while retaining the evidence needed for subsequent operations.}

The three head forms used in the view calculation record identity,
an inclusion from $\bot$, or an inclusion into $\top$.
Function and object conversions need forms that retain their component
evidence or templates. For our fragment, the possibilities are:
\begin{center}
\renewcommand{\arraystretch}{1.2}
\begin{tabularx}{\linewidth}{@{}lX@{}}
\toprule
Form & What its typing establishes after resolving definitions\\
\midrule
$\form{bot}$ & The source endpoint is $\bot$.\\
$\form{top}$ & The target endpoint is $\top$.\\
$\form{id}$ & The two endpoints agree.\\
$\form{pi}(d,c)$ & Both endpoints are functions; $d,c$ adapt argument and result.\\
$\form{obj}(E_0,\ldots,E_n)$ & The object conversion has these normalized recipes.\\
\bottomrule
\end{tabularx}
\end{center}
This table covers the forms used in the running examples. The full
target also has forms for self-bounds, discussed in
Section~\ref{sec:larger}.

The form $\form{id}$ carries no endpoints; its typing judgment requires
them to resolve equally. It is therefore a valid form for $\ell.A\le W$
when the store defines $\ell.A$ as $W$.

\subsection*{Normalizing the round trip}
The view at $\ty{Exact}$ gives
\[
 \begin{aligned}
 \hnf{\ev{lower}_A(a^\circ)}&=\form{id},\\
 \hnf{\ev{upper}_A(a^\circ)}&=\form{id}.
 \end{aligned}
\]
For composition the normalizer combines the two head forms,
using the operation on forms introduced in the view calculation:
\[
 \hnf{\ev{lower}_A(a^\circ);\ev{upper}_A(a^\circ)}
   =\form{id};\form{id}=\form{id}.
\]
The original inclusions have endpoints $W\le \ell.A$ and $\ell.A\le W$,
and their composite has endpoints $W\le W$. Resolution justifies the
identity forms at the intermediate endpoints.

\subsection*{Composing interface coercions}
The coercions $\ev{expose}$ and $\ev{forget}$ can also be composed
before either is applied to $\ell$.
The lower recipe of $\ev{forget}$ is
\[
 [\,\ev{bot}(W)\mid\text{inclusion slot }0\mid\text{absent}\,].
\]
Slot 0 of its input is the lower recipe produced by
$\ev{expose}$:
\[
 [\,\text{absent}\mid\text{equality slot }0\text{ reversed}\mid
       \text{absent}\,].
\]
Substituting the second recipe into the first hole yields
\[
 [\,\ev{bot}(W)\mid\text{equality slot }0\text{ reversed}\mid
       \text{absent}\,].
\]
The upper composite is
\[
 [\,\text{absent}\mid\text{equality slot }0\text{ forward}\mid
       \ev{top}(W)\,].
\]
The resulting morphism converts $\ty{Lit}$ directly to $\ty{Abs}$.
Normalizing its sides leaves the equality holes intact; instantiating
them with $\view{\ell}$ gives $\form{bot}$ and $\form{top}$, as with the
two successive view transformations. The morphism representation thus
supports composition by an explicit substitution operation on finite
templates.

\subsection*{Scope of head normalization}
The representation determines how far normalization proceeds.
Basic object templates access self through designated slots, with side
evidence outside the self binder. Composition can therefore combine
existing components without introducing a new member lookup through self.

A function form $\form{pi}(d,c)$ retains its argument and result
adaptations as evidence. In the dependent case, $c$ may mention a parameter
whose value has not yet been supplied. The machine needs only the outer
function form to perform a call; once the argument is available, it uses
$d$ and $c$ as casts. Head normalization therefore does not require
store-based normalization under an unapplied parameter.

The executable normalizer in Lean takes \emph{fuel}: a natural-number
budget that decreases during recursion. It terminates on every input
and returns an optional form, with failure possible both for malformed
evidence and for insufficient fuel. Failure at a particular budget
therefore does not establish that the evidence is ill typed.

The normalization theorem applies to evidence \emph{closed over a typed
store}: every free variable names a location in that store, whose value
has been checked at its precise type. Such evidence may mention $\ell$;
its checking context contains no extra opaque parameter assumptions.
For well-typed evidence in this setting, there is a fuel budget for which
normalization succeeds with a form of the required type. The proof is by
structural induction on typing derivations.
This is the inclusion-evidence part of the \emph{canonical-forms theorem}.

\exercise{compose}{Fill the slot in
$[\,\ev{bot}(W)\mid\text{equality 0 reversed}\mid\text{absent}\,]$
using the concrete view $[\,\form{eq}\,]$. Which head form results?
Why is the unfilled recipe not already a concrete view?}

\lesson{Function coercions and application}\label{sec:application}
\focus{A function head form supplies the argument and result adaptations required by an application.}

In the store constructed in Section~\ref{sec:closed-program}, location
$z$ holds $\lambda(u:\top).u$ at precise type $W=\top\to\top$.
Using the reopened exact-interface atom $a^\circ=R(a)$, that section's
closing calculation gives
\[
 z'=\cast{\cast{z}{\ev{lower}_A(a^\circ)}}{\ev{upper}_A(a^\circ)}:W.
\]
The chain of casts from $z$ to $z'$ normalizes to $\form{id}$.
If $r:\top$ is an argument atom, the machine can therefore call the
lambda stored at $z$ directly. Substitution in the body $u$ returns $r$.
For the original program, choose
$r=\cast{z}{\ev{top}(W)}$.
The result erases to $z$, as expected from evaluation of the identity.

Establishing this step requires more than the type $z':W$: the atom
$z'$ is a cast expression, while application must locate a stored lambda.
The canonical-forms theorem connects the atom's presented type to the
value at its root.

\subsection*{One machine step, with its premises}
Write $\langle\sigma,K,t\rangle$ for a state with store $\sigma$,
continuation $K$, and current term $t$. Let
$\operatorname{chain}_\sigma(c)$ normalize the wrappers from an atom's
root to its presented type. For a wrapped receiver, the identity case
of the machine's application rule is
\[
 \frac{\sigma(\ev{root}(c))=\lambda(u:S).t_0
       \qquad c\ne\ev{root}(c)
       \qquad \operatorname{chain}_\sigma(c)=\form{id}}
      {\langle\sigma,K,c\,r\rangle
         \longrightarrow\langle\sigma,K,t_0[u:=r]\rangle}.
\]
The comparison in the middle treats the root as a bare atom.
The chain equation abbreviates successful normalization at some fuel;
it is a premise of the step, not an unchecked optimization.

For the closed program, set $c=z'$, $K=\epsilon$, $S=\top$, and
$t_0=u$. The root is $z$, its stored value is the identity lambda,
and both member casts have identity form by the preceding view
calculation. Hence all three premises hold, and the actual step is
\[
 \langle\sigma,\epsilon,z'\,\cast{z}{\ev{top}(W)}\rangle
 \longrightarrow
 \langle\sigma,\epsilon,\cast{z}{\ev{top}(W)}\rangle.
\]
This is \texttt{Step.appCastRefl}. It changes neither the store nor the
continuation, and its result is an atom. Earlier allocation and alias
substitution took several steps; this display isolates one complete
transition and all the evidence needed to enable it.

\subsection*{Adapting an argument}
For a nonidentity function coercion, present $z$ at the function type
$W\to\top$, whose domain is narrower than that of $W$.
The required inclusion is
\[
  (\top\to\top)\le(W\to\top).
\]
Clients of the new interface supply arguments at $W$, which the original
function can accept at $\top$. The argument and result evidence are
\[
 d=\ev{top}(W):W\le\top,\qquad
 c=\ev{refl}(\top):\top\le\top.
\]
The function constructor builds
$\ev{pi}(d,c):(\top\to\top)\le(W\to\top)$.
Head normalization exposes its function form:
\[
 \hnf{\ev{pi}(d,c)}=\form{pi}(d,c).
\]
The components $d,c$ remain evidence for the call's argument and result
casts. The domain adaptation is contravariant: the function inclusion
maps the original function interface to the client interface, while
$d$ maps the client's argument type to the original domain.

Figure~\ref{fig:function-call} gives the source rule, its target
evidence, and the call. A client supplies $z:W$; the argument evidence
presents it at $\top$ for the stored identity. The body returns that
same atom, and the result evidence presents it to the client.
Administrative steps arrange the pending casts around evaluation of the
function body. The call erases to $z\,z$ and its result to $z$.


\begin{figure}[!htbp]
\begin{examplepanel}{srccol}{Source function subtyping --- contravariant domain}
\[
\src{\inferrule*[right=\textsc{All}]{
 W<:\top\\ \top<:\top
}{(\top\to\top)<:(W\to\top)}}
\]
\end{examplepanel}
\begin{examplepanel}{tgtcol}{Target evidence and application}
\[
\tgt{\begin{aligned}
 d&=\ev{top}(W):W\le\top,&c&=\ev{refl}(\top):\top\le\top,\\
 g&=\cast{z}{\ev{pi}(d,c)}:W\to\top.
\end{aligned}}
\]
\par\small
\renewcommand{\arraystretch}{1.2}
\noindent\begin{tabularx}{\linewidth}{@{}lX@{}}
\toprule
Stage of \(g\,z\) & Argument or result at that stage\\
\midrule
Client supplies & $z:W$\\
Original function receives & $\cast{z}{d}:\top$\\
Identity body returns & $\cast{z}{d}:\top$\\
Client receives & $\cast{\cast{z}{d}}{c}:\top$\\
\bottomrule
\end{tabularx}\par
\end{examplepanel}
\begin{examplepanel}{runcol}{Erased execution}
\[
\run{\erase{\sigma}(z)=\lambda u.u,\qquad z\,z\longrightarrow z.}
\]
\end{examplepanel}
\caption{The function coercion adapts the argument with \(d\) and the
result with \(c\). The typed store holds \(z=\lambda(u:\top).u\), and
the function head form retains both evidence components for application.
The table groups administrative steps; the erased call applies the
stored identity.}
\label{fig:function-call}
\end{figure}



\FloatBarrier
\needspace{8\baselineskip}
\subsection*{Dependent results and composition}
For nondependent function types, $c$ is independent of the argument.
For a dependent function $(u:S)\to T(u)$, the result type can mention
$u.A$. At a call, substitute the argument atom in term positions and its
root in type positions. A wrapper changing the argument's interface
must not create a different dependent result name.

Composition of function coercions respects this order of adaptation.
If a first adaptation expects an intermediate input and a second offers
a still narrower one, the second input adaptation runs first.
Result adaptations run in the order of the function coercions.
In a dependent codomain, the earlier result evidence is instantiated
with the adapted argument before being composed with the later result evidence.

\remember{Application uses a function head form to adapt the argument,
evaluate the stored lambda, and adapt the result. The component evidence
is retained until the argument is available.}

\exercise{variance}{Suppose $d:S_2\le S_1$ and $c:T_1\le T_2$,
with nondependent results. Give the function inclusion built from them.
When it is applied to $b:S_2$, at which type does the original body
receive $b$, and at which type does the client receive the result?}

\lesson{Bad bounds and uninhabited interfaces}\label{sec:bad}
\focus{Inconsistent bounds can occur in an open context, but no typed store supplies an argument satisfying them.}

Consider the source interface $\iface{\top}{\bot}$ and its target
translation
\[
 \ty{Bad}=\obj{\top\le s.A,\ s.A\le\bot}.
\]
Under a hypothetical parameter $y:\ty{Bad}$, we can read both facts
and compose them:
\[
 \ev{lower}_A(y);\ev{upper}_A(y):\top\le\bot.
\]
This evidence is well typed in the open context. In particular, a
function body may be checked under an uninhabited parameter type.
The safety argument must show that execution cannot supply an argument
at that type.

\needspace{10\baselineskip}
\subsection*{The obstruction to constructing an argument}
Our stored object defines $A=W$.
To present it at $\ty{Bad}$ using the recipe pattern, we would need
both chains:
\[
 \begin{aligned}
 \text{lower: }&\top\le W\le \ell.A,\\
 \text{upper: }&\ell.A\le W\le\bot.
 \end{aligned}
\]
The definition equality handles the links next to $\ell.A$.
The remaining obligations are $\top\le W$ and $W\le\bot$.
Composing a successful construction would produce closed evidence for
$\top\le\bot$ over the typed store.

The normalization theorem excludes such evidence. None of the typed
head forms for this fragment admits those endpoints:
$\top$ is not $\bot$; $\bot$ is not $\top$; they do not resolve equally;
they are not both functions; and they are not both object interfaces.
The full target's consistency result covers its additional forms as well.

The same consistency result excludes every atom at $\ty{Bad}$ over a
typed store: selecting and composing its bounds would again yield
closed evidence for $\top\le\bot$.

\subsection*{Open assumptions and store contexts}
The open derivation used $y:\ty{Bad}$ as an assumption.
The runtime theorem requires a store context, which records the stored
values' precise types and definitions. Well-scopedness or well-formedness of
an open context alone does not establish this hypothesis.

At application, substitution connects the open body, checked under its
opaque parameter, to a typed atom over the store. No such atom exists
for this parameter type, so execution cannot reach the body through a
well-typed call. The argument does not require normalizing the open
evidence under the impossible assumption.

\exercise{bad}{Consider the claim that the well-typed open evidence
$\ev{lower}_A(y);\ev{upper}_A(y)$ must have a typed head form, proving
the system inconsistent. Identify the missing theorem hypothesis precisely.}

\lesson{Canonical forms and type safety}\label{sec:proof}
\focus{The canonical-forms theorem connects evidence normalization to the operations of the target machine.}

The proof establishes properties of several judgments simultaneously,
following the dependencies in the syntax: cast atoms contain inclusion evidence, member selection
requires an atom, and object coercions contain templates. Normalization
of inclusion evidence therefore depends on properties of atoms and views.

\needspace{15\baselineskip}
\subsection*{The canonical-forms theorem}
Fix a typed store $\sigma$ with its context $\Gamma_\sigma$.
All judgments below are in that store context, so their free variables
are allocated locations. The canonical-forms theorem establishes the
following related facts:
\begin{enumerate}
\item Well-typed inclusion evidence has a typed head form.
\item Equality evidence relates endpoints that resolve equally.
\item A well-typed atom has a typed view at its object interface.
\item The atom's chain of wrappers has a form relating the root's
precise type to its presented type.
\item For the full target, field-presence evidence identifies a field actually stored in
the appropriate object.
\end{enumerate}
The proof also excludes atoms whose types resolve to $\bot$. Alongside
successful normalization, these statements establish the typing and
semantic properties of the resulting forms.

Three representative cases illustrate the dependencies. For transitivity,
the induction hypotheses
provide typed forms for both components; a composition lemma shows that
combining them preserves typing. For a member read, the hypotheses provide the atom's
view and the interface conversion's normalized recipes; selecting the
appropriate resulting entry gives the needed fact. For a cast atom,
the proof transforms its existing view and composes its chain.
Each case uses an operation from the preceding calculations.

\needspace{8\baselineskip}
\subsection*{Target progress and preservation}
A machine state consists of a store, a current expression, and a
continuation recording pending lets and casts.

\emph{Progress} says a well-typed machine state is final or can step.
At application, the canonical-forms facts identify a lambda at the
function atom's root and give the required argument and result evidence.
At field projection, presence evidence identifies a field to read.
The remaining cases handle lets and casts.

\emph{Preservation} says a step keeps the program and store typed,
allowing for newly allocated locations. Substitution must transport
evidence as well as terms. Allocation must keep the bare value at its
precise type and move its presentations into the body. A function call
must cast the argument and result at the endpoints promised by its form.

Together, progress and preservation ensure that a closed well-typed
target program never reaches a stuck state. The program may diverge;
termination of evidence normalization does not imply termination of
program execution.

\subsection*{Transporting safety to the source}
The translation takes a source \emph{typing derivation}.
For a well-formed source context, it produces a well-typed target term.
Different source derivations can insert different evidence,
but the translation erases to the same underlying source program:
\[
 \erase{\operatorname{translate}(h)}=\erase{t}
 \quad\text{when }h\text{ derives }\Gamma\vdash t:T.
\]
Transporting safety additionally requires operational correspondence.

Both source and target erase into a shared runtime. A simulated source
state is one for which a typed target state has the same erasure.
Translation establishes this relation initially. When the source steps,
runtime correspondence and target simulation provide a matching target
execution; preservation re-establishes the relation.
Conversely, target progress, together with reflection of runtime steps
by the source, shows the source state is final or can step.
Here \emph{reflection} means that a runtime step out of an erased source
state comes from a corresponding source step. Pending administrative
cast steps can be cleared in finitely many steps before making this
argument. Thus no infinite execution can consist solely of administrative
steps. These additional target steps explain
why the two machines do not run in lockstep.

\remember{Canonical forms establish the runtime structure needed for
target safety. Translation and operational correspondence then transport
safety to the source: erasure relates program syntax, and simulation
relates executions.}

For the main source calculus, this argument derives safety from the
target's metatheory and the simulations, without a separate source
preservation proof as a premise.

\exercise{proof}{Which fact is missing from each proposed proof?
(a) A well-typed translation is claimed to establish source safety by itself.
(b) The existence of one target step is claimed to establish safety for
all future steps. Give a different missing ingredient for each.}
